\section{Introduction}

Weak-form market efficiency remains a core benchmark in empirical finance. In emerging markets, however, integration frictions, episodic illiquidity, and structural regime changes motivate repeated reassessment of short-horizon predictability \parencite{fama1970,fama1991,harvey1995,bekaertharvey1995}.

This study asks one focused question:
\begin{quote}
Do diffusion-based forecasts deliver robust out-of-sample loss reductions relative to strong benchmark models for Latin American ETF cumulative returns?
\end{quote}

The empirical contribution is a benchmark-disciplined comparison across EWZ, EWW, ECH, and GXG over six $(L,H)$ configurations with expanding-window evaluation and Diebold-Mariano inference against random walk. The methodological contribution is a leakage-audited, horizon-consistent workflow that combines classical baselines, nonlinear baselines, conditional diffusion, and seed/feature sensitivity analysis under a common target definition. The interpretability contribution is a stability-oriented XAI layer that treats feature reliance as predictive model dependence, not causal mechanism.

Inference is intentionally conservative. The design supports statements about forecast association and incremental predictive content; it does not identify causal effects, structural transmission parameters, or implementable profitability after transaction costs.
